
\documentclass[handout]{beamer}

%\usepackage[table]{xcolor}
\mode<presentation> {
  \usetheme{Boadilla}

%  \usetheme{Pittsburgh}
%\usefonttheme[2]{sans}
\renewcommand{\familydefault}{cmss}
%\usepackage{lmodern}
%\usepackage[T1]{fontenc}
%\usepackage{palatino}
%\usepackage{cmbright}
  \setbeamercovered{transparent}
\useinnertheme{rectangles}
}

\setbeamercolor{normal text}{fg=black}
\setbeamercolor{structure}{fg= blue}
\definecolor{trial}{cmyk}{1,0,0, 0}
\definecolor{trial2}{cmyk}{0.00,0,1, 0}
\definecolor{darkgreen}{rgb}{0,.4, 0.1}
\usepackage{array}
\beamertemplatesolidbackgroundcolor{white}  \setbeamercolor{alerted
text}{fg=red}
\usepackage{tcolorbox}
\setbeamertemplate{caption}[numbered]\newcounter{mylastframe}


\usepackage{tikz}
\usetikzlibrary{arrows}
\usepackage{colortbl}

\renewcommand{\familydefault}{cmss}
%\usepackage[all]{xy}

\usepackage{tikz}
\usepackage{lipsum}

 \newenvironment{changemargin}[3]{%
 \begin{list}{}{%
 \setlength{\topsep}{0pt}%
 \setlength{\leftmargin}{#1}%
 \setlength{\rightmargin}{#2}%
 \setlength{\topmargin}{#3}%
 \setlength{\listparindent}{\parindent}%
 \setlength{\itemindent}{\parindent}%
 \setlength{\parsep}{\parskip}%
 }%
\item[]}{\end{list}}
\usetikzlibrary{arrows}

\usecolortheme{lily}

\newtheorem{com}{Comment}
\newtheorem{lem} {Lemma}
\newtheorem{prop}{Proposition}
\newtheorem{condition}{Condition}
\newtheorem{thm}{Theorem}
\newtheorem{defn}{Definition}
\newtheorem{cor}{Corollary}
\newtheorem{obs}{Observation}
 \numberwithin{equation}{section}
 

\setbeamertemplate{navigation symbols}{}


\title[PLSC 43401]{Mathematical Foundations of Political Methodology}

\author{Robert Gulotty}
\institute[Chicago]{University of Chicago}
\vspace{0.3in}




\begin{document}

\begin{frame}
\maketitle
\end{frame}

\begin{frame}
\frametitle{Day 2 : Conditional probability}

\begin{itemize}
\item Conditional Probability
\item Bayes Rule
\item Independence
\item Counting
\end{itemize}
\end{frame}

\begin{frame}{Why study conditional probability?}

Most of the relationships we are interested in as policy makers/analysts, are conditional ones. Knowing that an event has occurred, might lead us to reassess the probability of another event.\pause

\begin{itemize}
\item What is the probability of Iran obtaining nuclear weapons given that it has a uranium enrichment program?\pause
\item What is the probability that we will be able to affect climate change if we adopt the climate change treaty?\pause
\item What is the probability of holding liberal political views given certain demographic characteristics?
\end{itemize}
\end{frame}

\begin{frame}
\frametitle{Definition of Conditional probability}
The \emph{conditional probability} of an event A, given an event B is given by:
$$P(A|B)=\frac{P(A\cap B)}{P(B)}$$
for all $P(B)\neq 0$\\
\pause We can rearrange this to form the``Multiplication Rule'':
$$P(A\cap B)=P(A|B)\times P(B)$$

\end{frame}



\begin{frame}
\frametitle{Examples of Conditional Probability} 

Example 1: \pause  
\begin{itemize}
\item[-] $F = \{\text{Myanmar sells oil to China} \} $\\  \pause 
\item[-] $E = \{\text{Myanmar discovers oil} \} $ \\  \pause 
\item[-] If $F$ occurs then $E$ must occur, $P(E|F) = 1$ \\  \pause 
\end{itemize}


Example 2: Security Council voting on Syria strike \pause 
\begin{itemize}
\item[-] $V = \{ \text{Russia and China veto strike}\} $\pause
\item[-] $A = \{ \text{Strike approved by UN}\}$  \pause 
\item[-] $P(A|V)  = \frac{P(A \cap V)}{P(V) } = \frac{P(\emptyset)}{P(V)}$ = 0
\end{itemize}


\end{frame}




\begin{frame}
\frametitle{Axioms of Conditional probability}
$P(A|B)$ is a probability, so it follows the axioms.
\begin{enumerate}
\item $P(A|B)\geq 0, \ \ \forall A$
\item $P(\Omega|B)=1$.
\item For all mutually exclusive events, $A_1,\ A_2,\ ... A_n$
$$P(A_1\cup A_2 \cup...\cup A_n|B)= P(A_1|B)+P(A_2|B)+...+ P(A_n|B)$$
\end{enumerate}
\end{frame}

\begin{frame}
\frametitle{Conditional Probability Example}
Basic claim: Foreign military threats (given certain other conditions) produce revolutions.
\begin{itemize}
\item Skocpol (1979)
\item Let E denote the event of a revolution and
\item Let F the presence of a foreign threat 
\end{itemize}
The way to read Skocpol's hypothesis in probability language is:\\
$P(E|F)\geq P(E|F^c)$

\end{frame}


\begin{frame}
\frametitle{Conditional Probability Example}
How might we know whether $P(E|F)>P(E|F^c)?$\pause
\begin{itemize}
\item $P(E|F)=\frac{P(E\cap F)}{P(F)}$
\item $P(E|F^c)=\frac{P(E\cap F^c)}{P(F^c)}$
\item So we need to know both the rate of Revolutions with foreign threats, the rate of Revolutions without foreign threats, and the probability of a foreign threat.
\end{itemize}

\end{frame}

\begin{frame}
\frametitle{Revolutions}
Skocpol's data:\\
\vspace{1 cm}

\begin{tabular}{r|cc}
&Revolution&No Revolution\\
\hline
Foreign Threat& France, Russia, China\\
No Threat&& England, Prussia, Japan\\
\end{tabular}
\end{frame}



\begin{frame}
\frametitle{$P(A|B)$ and $P(B|A)$ are usually different}


\begin{eqnarray}
P(A|B) & = & \frac{P(A\cap B)}{P(B)} \nonumber \\ 
P(B|A) & = & \frac{P(A \cap B) } {P(A)} \nonumber  
\end{eqnarray}


\pause \pause \pause 
\textbf{Example}: Race and Welfare\\ \pause 
`I will go to the NAACP convention and tell the African-American community why they should demand paychecks instead of food stamps.''\\
			\hfill -Newt Gingrich, January 5, 2012 \pause 

\vspace{0.25in}
Consider new enrollments to TANF \pause 
\begin{itemize}
\item[] P(TANF Recipient$|$ Black )  $\approx$ 29 \%\\  \pause 
P(TANF Recipient$|$ White )  $\approx$ 11\% \\  \pause 
\item[] P(Black$|$TANF Recipient ) $\approx$ 26 \% \pause 
\item[] P(White$|$TANF Recipient ) $\approx$ 47 \% \pause 
\end{itemize}

Confusing $P(A|B)$ and $P(B|A)$ leads to incorrect inference about welfare.

\end{frame}


\begin{frame}
\frametitle{Law of Total Probability}

For a set of mutually exclusive events, $E_1,\  E_2,\ E_3 .... E_n$, which together form the sample space $\Omega$, the probability of any event A is:

$$P(A)= P(A|E_1)P(E_1)+P(A|E_2)P(E_2)+...+P(A|E_n)P(E_n)$$
\pause Why?\pause
$$P(A)= P(A\cap \Omega)$$\pause
$$P(A)= P(A\cap (E_1\cup E_2 \cup ... \cup E_n))$$\pause
$$P(A)= P((A\cap E_1)\cup (A\cap E_2) \cup ... \cup(A\cap E_n))$$\pause
$$P(A)= P(A\cap E_1)+ P(A\cap E_2) + ...P(A\cap E_n)$$\pause
$$P(A)= P(A|E_1)P(E_1)+P(A|E_2)P(E_2)+...+P(A|E_n)P(E_n)$$

\end{frame}


\begin{frame}
\frametitle{Example of Law of Total Probability}
 Say you decide to implement a Democracy Promotion Program (DPP) in some villages in Myanmar in the hopes of increasing voter turnout at the next election. What is P(vote) after DPP?\pause
 \begin{itemize}
 \item Let $P(vote|DPP)=0.6$\pause
 \item Let $P(vote|\sim DPP)=0.3$\pause
 \item Let $P(DPP) = P(\sim DPP) = 0.5$ (Decided based on coin-flipping)\pause
 \item $P(vote)= .6*.5+.3*.5= .45$
 \end{itemize}
\end{frame}


\begin{frame}
\frametitle{Example: Law of Total Probability} 

Infer $P(\text{vote})$ after mobilization campaign \pause 
\begin{itemize}
\item[-] $P(\text{vote}|\text{mobilized} ) = 0.75$ \pause 
\item[-] $P(\text{vote}| \text{not mobilized} ) = 0.25 $ \pause 
\item[-] $P(\text{mobilized}) = 0.6 ; P(\text{not mobilized} ) = 0.4 $ \pause 
\item[-] What is $P(\text{vote})$? \pause 
\end{itemize}

Sample space (one person) = \\$\{$ (mobilized, vote), (mobilized, not vote), (not mobilized, vote) , (not mobilized, not vote) $\}$\\
\alert{Mobilization} partitions the space (mutually exclusive and exhaustive)  \pause \\
We can use the law of total probability \pause 
\begin{eqnarray}
P(\text{vote} ) & = & P(\text{mob.} ) \times P(\text{vote}| \text{mob.} ) + P(\text{not mob}) \times P(\text{vote} | \text{not mob} ) \nonumber \\
 & = & 0.6 \times 0.75 + 0.4 \times 0.25 \nonumber \\
  & = & 0.55 \nonumber  
  \end{eqnarray}

\end{frame}


\begin{frame}
\frametitle{Example Law of Total Probability}

\alert{Mixture Models}: flexible modeling strategy.  \\ \pause 

Two coins:\\ \pause 
\begin{itemize}
\item[-] Fair: $P(H) = 1/2$  \pause
\item[-] Biased $P(H) = 3/4$ \pause 
\end{itemize}

Draw a coin from urn ($P(\text{fair}) = 1/2)$ and then flip.  \\
$P(H)?$ \pause 

$\Omega = \{(\text{fair}, H), (\text{fair}, T), (\text{bias}, H), (\text{bias}, T) \}$ \pause 

\begin{eqnarray}
P(H) & = & P(\text{fair}) \times P(H| \text{fair} )  + P(\text{bias} ) \times P(H | \text{bias} ) \nonumber\\ \pause 
& = & \frac{1}{2} \times \frac{1}{2} + \frac{1}{2} \times \frac{3}{4} \nonumber \\
& = & \frac{5}{8}\nonumber \pause 
\end{eqnarray}

\alert{Mixture} of two coins


\end{frame}






\begin{frame}
\frametitle{Bayes' Rule} 

\begin{itemize}
\item[-] $P(B|A)$ may be easy to obtain
\item[-] $P(A|B)$ may be harder to determine
\item[-] Bayes' rule provides a method to move from $P(B|A)$ to $P(A|B)$.  
\end{itemize}



\end{frame}

\begin{frame}
\frametitle{Conditional Probability and Bayes' Rule}
 $$P(A|B)=\frac{P(A\cap B)}{P(B)}$$
 $$P(A\cap B)=P(B\cap A)=P(B|A)P(A)$$\pause
 
 \begin{thm}{Bayes' Rule}
  $$P(A|B)=\frac{P(B|A)P(A)}{P(B)}$$
\end{thm}
\end{frame}


\begin{frame}
\frametitle{General Statement of Bayes' Rule}
Let the events $A_1...A_k$ form a partition of the space S such that $P(A_j)>0$ for $j=1...k$.
And let B be an event such that $P(B)>0$.  Then for $i=1...k$
 
 \begin{thm}{Bayes' Rule}
  $$P(A_i|B)=\frac{P(A_i)P(B|A_i)}{\sum_{j=1}^{k}P(A_j)P(B|A_j)}$$
\end{thm}
\end{frame}




\begin{frame}
\frametitle{Bayes' Rule: Example} 

How do we identify racial groups from lists of names?  \pause \\
Census Bureau collects information on distribution of names by race. \pause  \\
For example, \alert{Washington} is the ``\alert{blackest}" name in America. \pause\\
\begin{itemize}
\item[-] P(black)= 0.126.
\item[-] P(not black) = 1 - P(black) = 0.874. 
\item[-] P(Washington$|$ black) = 0.00378.
\item[-] P(Washington$|$ not black) = 0.000060615. \pause 
\end{itemize}

\begin{eqnarray}
P(\text{black}|\text{Wash} ) & = & \frac{P(\text{black}) P(\text{Wash}| \text{black}) }{P(\text{Wash} ) } \nonumber \\
 & = & \frac{P(\text{black}) P(\text{Wash}| \text{black}) }{P(\text{black})P(\text{Wash}|\text{black}) + P(\text{nb})P(\text{Wash}| \text{nb}) } \nonumber\\
 & = & \frac{0.126 \times 0.00378}{0.126\times 0.00378 + 0.874 \times 0.000060616} \nonumber\\
 & \approx & 0.9 \nonumber  
 \end{eqnarray}


\end{frame}




\begin{frame}
\frametitle{Monty Hall Problem}

Suppose we have three doors.  $A, B, C$.  \pause \\
Suppose you're on a game show, and you're given the choice of three doors: Behind one door is a car; behind the others, goats. \\ \pause 
\begin{itemize}
\item[-] A contestant guesses a door. \pause 
\item[-] The host opens a different door, revealing one of two goats. The contestant has option to switch. \pause 
\item[-] Should the contestant switch? \pause 
\end{itemize}

Contestant guesses $A$\\ \pause 
$P(A) = 1/3 \leadsto$ chance of winning without switch\\\pause 
If $C$ is revealed to not have a car:\\ \pause 



\end{frame}


\begin{frame}
\frametitle{Monty Hall Problem}

\pause 
\begin{eqnarray}
P(B| C \text{ revealed} ) & = & \frac{P(B)P(C \text{ revealed} | B)}{P(B)P(C \text{ revealed} | B) + P(A) P(C \text{ revealed} | A) } \nonumber \\ \pause 
& = & \frac{1/3 \times 1}{1/3 + 1/3 \times 1/2 } = \frac{1/3}{1/2} = \frac{2}{3}\nonumber \\ \pause 
P(A| C \text{ revealed} ) & = & \frac{P(A) P(C \text{ revealed} | A)}{ P(B)P(C \text{ revealed} | B) + P(A) P(C \text{ revealed} | A) }\nonumber \\ \pause 
& = & \frac{1/3 \times 1/2}{1/3 + 1/3 \times 1/2} = \frac{1}{3} \nonumber  \pause 
\end{eqnarray}

\alert{Double chances of winning with switch}


\end{frame}





\begin{frame}
\frametitle{Independence and Information} 

Does one event provide \alert{information} about another event? \pause 

\begin{defn} 
If learning about B does not change the probability that A occurs, A is independent of B.
 $$P(A|B)=P(A)$$\pause
 $$P(A|B)=\frac{P(A\cap B)}{P(B)}=P(A)$$
\end{defn}
\pause

\invisible<1-2>{Independence is symmetric: if $F$ is independent of $E$, then
  $E$ is independent of $F$} 

\end{frame}




\begin{frame}
\frametitle{Independence}
What if A is \emph{Mutually Exclusive} of B?\pause
 $$A\cap B=\emptyset$$\pause
 
  $$P(A|B)=\frac{P(\emptyset)}{P(B)}=0\neq P(A)$$
 \pause
 If two events occur with positive probability, then they cannot be mutually exclusive and independent.
 \end{frame}




\begin{frame}
\frametitle{Example Independence Relationship}

Flip a fair coin twice.  
\begin{itemize}
\item[] $E = \text{first flip heads} $
\item[] $F = \text{second flip heads} $
\end{itemize}\pause
\begin{eqnarray}
P(E \cap F ) & = & P( \{ (H, H) , (H, T) \} \cap \{ (H, H), (T, H) \} ) \nonumber \pause \\
 & =& P( \{(H, H)\} ) \nonumber \\ 
 & = & \frac{1}{4} \nonumber \\
P(E ) & = & \frac{1} {2} \nonumber \pause \\
P(F) & = & \frac{1}{2} \nonumber \pause \\
P(E)P(F)  & =& \frac{1}{2}\frac{1}{2} = \frac{1}{4}  =P(E \cap F )  \nonumber 
 \end{eqnarray} 
 
 

\end{frame}

\begin{frame}
\frametitle{Independence: No Information} 

Suppose $E$ and $F$ are independent.  Then, 
\begin{eqnarray}
P(E|F ) & = & \frac{P(E \cap F) }{P(F) } \nonumber \\
& = & \frac{P(E)P(F)}{P(F)}\nonumber \\
& = & P(E) \nonumber 
\end{eqnarray}


Conditioning on the event $F$ does not modify the probability of $E$.  \\
\alert{No information about $E$ in F}  \\

\end{frame}



\begin{frame}
\frametitle{Independence and Compliments}

\begin{prop}
Suppose $A$ and $B$ are independent events. Then the events $A$ and $B^{c}$ are also independent.  
\end{prop}

\pause 

\begin{proof}
\begin{eqnarray}
P(A \cap B^{c} ) & = & P(A) - P(A \cap B) \nonumber  \pause \\
& = & P(A)  - P(A) P(B) \nonumber  \pause \\
& = & P(A) ( 1- P(B) )  \pause \nonumber \\
& = & P(A) P(B^{c})   \nonumber 
\end{eqnarray}
\end{proof}



\end{frame}


\begin{frame}
\frametitle{Independence and Causal Inference}
 We want to know what the effect of some treatment is on some outcome of interest.\pause
 \begin{itemize}
\item Democracy promotion programs on democracy \pause
\item Cash transfers on education or health outcomes \pause
\end{itemize}
In most studies, inference is made based on observations: observational studies. Units select into treatment, or are selected. \pause
\end{frame}




\begin{frame}{Example of Independence and Causal inference.}
Example: democracy promotion programs in countries where we think they are likely to succeed\pause

\begin{table}
\begin{tabular}{lll}
		&Democracy& No Democracy\\
DPP &10 & 50\\
no DPP & 2& 13\\

\end{tabular}
\end{table}
\pause
$$P(Democracy|\ DPP)= \frac{(10/75)}{(60/75)} =1/6= .17$$
$$P(Democracy| no\ DPP)=  \frac{(2 /75)}{(15/75)} = 2/15=.13$$

 \end{frame}




\begin{frame}{Example of Independence and Causal inference.}
Example: democracy promotion programs in countries where we think they are likely to succeed\pause

\begin{table}
\begin{tabular}{llll}
		&&Democracy& No Democracy\\
 Good Gov & DPP &2 & 18\\
		 &no DPP & 1& 9\\
		 \hline
Bad Gov  & DPP &8 & 32\\
		&no DPP & 1& 4\\



\end{tabular}
\end{table}
\pause
$$P(Democracy| DPP\& GG)= \frac{(2/75)}{( 20/75)} = .1$$
$$P(Democracy| no DPP\& GG)=  \frac{(1/75)}{(10/75)} = .1$$

 \end{frame}




\begin{frame}
\frametitle{Conditional  Independence}

\begin{defn}
Let $E_{1}$ and $E_{2}$ be two events.  We will say that the events are conditionally independent given $E_{3}$ if 
\begin{eqnarray}
P(E_{1} \cap E_{2} | \alert{E_{3}}) & = & P(E_{1} | E_{3} ) P(E_{2} | E_{3} ) \nonumber 
\end{eqnarray}


\end{defn}



\end{frame}




\begin{frame}
\frametitle{Example: Independence and Causal Inference} 

Selection and Observational Studies
\begin{itemize}
\item[-] We often want to infer the effect of some treatment 
\begin{itemize}
\item[-] Incumbency on vote return 
\item[-] Democracy on war
\end{itemize}
\item[-] Observational studies: observe what we see to make inference 
\item[-] Problem: units select into treatment 
\begin{itemize}
\item[-] Simple example: enroll in job training if I think it will help 
\item[-] P(job$|$training in study) $\neq$ P(job$|$forced training) 
\end{itemize}
\item[-] \alert{Background characteristic}: difference between
  treatment and control groups
\item[-] \alert{Experiments}: make background characteristics and treatment status
  independent 
\end{itemize}


\end{frame}


\begin{frame}
\frametitle{Conditional Probability}

\begin{defn}
Let $E_{1}$ and $E_{2}$ be two events.  We will say that the events are conditionally independent given $E_{3}$ if 
\begin{eqnarray}
P(E_{1} \cap E_{2} | \alert{E_{3}}) & = & P(E_{1} | E_{3} ) P(E_{2} | E_{3} ) \nonumber 
\end{eqnarray}


\end{defn}



\end{frame}


\begin{frame}

\begin{prop}
Suppose $E_{1}$ and $E_{2}$ and $E_{3}$ are events such that $P(E_{1} \cap E_{2})>0$ and $P(E_{2} \cap E_{3}) >0$.  Then $E_{1}$ and $E_{2}$ are conditionally independent given $E_{3} $ if and only if $P(E_{1} | E_{2} \cap E_{3})  = P(E_{1} | E_{3}) $.  
\end{prop}

\begin{proof}
Suppose $E_{1}$ and $E_{2}$ are conditionally independent given $E_{3}$.  Then 
\begin{eqnarray}
P(E_{1} \cap E_{2} |  E_{3}) & = & \frac{P(E_{1} \cap E_{2} \cap E_{3}  ) }{P(E_{3} )} \nonumber \\
& = & \frac{P(E_{3})P(E_{2} | E_{3} )P(E_{1} | E_{2} \cap E_{3} ) }{ P(E_{3} )}\nonumber \\
P(E_{1}|E_{3}) P(E_{2} | E_{3} ) & = & P(E_{2} | E_{3} )P(E_{1} | E_{2} \cap E_{3} )\nonumber \\
P(E_{1}|E_{3}) & = & P(E_{1} | E_{2} \cap E_{3} ) \nonumber 
\end{eqnarray}

\end{proof}





\end{frame}

\begin{frame}

\begin{proof}
Suppose $P(E_{1} | E_{2} \cap E_{3})  = P(E_{1} | E_{3}) $

\begin{eqnarray}
P(E_{1} \cap E_{2}| E_{3 } ) & = & P(E_{2} | E_{3} ) P(E_{1} | E_{2}\cap E_{3} ) \nonumber \\
 & = & P(E_{2} | E_{3} ) P(E_{1} | E_{3} ) \nonumber 
\end{eqnarray}

\end{proof}

\end{frame}








\begin{frame}

\begin{defn}
Suppose we have a sequence of events $E_{1}, E_{2}, \hdots, E_{n}$.  We say the sequence of events is \alert{mutually independent} if for each subset of the sequence,  $E_{i_{1}}, E_{i_{2}}, \hdots, E_{i_{j}}$ 
\begin{eqnarray}
P(E_{i_{1}}\cap E_{i_{2}}\cap \hdots \cap E_{i_{j}}) & = & \prod_{m = 1}^{j} P(E_{i_{m}}) \nonumber 
\end{eqnarray}
\end{defn}

\alert{For a sequence to be independent, every subset is independent }


\end{frame}


\begin{frame}
\frametitle{Most Different research design (Method of Agreement)}

Given a set of possible causes, A, B, C, D, and E, if we observe an outcome $x$ after A, B and C in one case and after A, D, and E in another, than we have evidence that A causes x.\pause
This is sometimes called the method of elimination, where we rule out competing explanations.\\
\pause
\begin{itemize}
\item We must assume that the cause is deterministic.
\item We must assume that one of these causes must be the cause (we aren't missing anything).
\item We have ruled out B, C, D and E as necessary conditions.
\end{itemize}

\end{frame}


\begin{frame}
\frametitle{Most Similar research design (Method of Difference)}

Given the complete set of possible causes, A, B, and C, if we observe an outcome $x$ after A, B and C and not after B and C alone, then we have evidence that A causes x.\pause
\begin{itemize}
\item We must assume that the cause is deterministic.
\item We must assume that all conditions other than A are exactly the same.
\item We have established that A is a necessary condition.
\end{itemize}


\end{frame}


\begin{frame}{A note on Mill's Method}
``Nothing can be more ludicrous than the sort of parodies on experimental reasoning which one is accustomed to meet with, not in popular discussion only, but in grave treatises, when the affairs of nations are the theme. [...] ``How can such or such causes have contributed to the prosperity of one country, when another has prospered without them?'' Whoever makes use of an argument of this kind, not intending to deceive, should be sent back to learn the elements of some one of the more easy physical sciences.''\\
\hfill-John Stuart Mill. 1872. A System of Logic
\end{frame}


\begin{frame}
\frametitle{A note on Mill's Method}
John Stewart Mill rejected the use of his method for social science, particularly comparative and international relations, but they are still influential.

\end{frame}






\end{document}